problem:  
Let $\phi:(M^{m},g)\to (N^{2},h)$ be a harmonic morphism whose fibers are compact and minimal.  
For a basic function $u:M\to\mathbb{R}$ (i.e., $u$ constant along fibers), the conformally related metric $\tilde g=e^{2u}g$ preserves the horizontally conformal submersion property of $\phi$.  
Let $\mathcal{H}$ denote the horizontal distribution, and let $A$ and $T$ denote its O’Neill tensors.  
Given a horizontal 2‑surface $\Sigma$ that is a local section ($\phi|_\Sigma:\Sigma\to N$ is a local diffeomorphism) and minimal with respect to $g$, write $\tau_u(\phi|_\Sigma)$ for the tension field of the map $\phi|_\Sigma:(\Sigma,\tilde g|_\Sigma)\to (N,h)$.  

**Problem:**  
Determine the conditions on $(M,g,\phi)$ under which there exists a nontrivial basic conformal factor $u$ satisfying  
\[
\tau_u(\phi|_\Sigma)=0 \quad \text{for every horizontal minimal section } \Sigma.
\]  
Equivalently, when can one conformally modify the total space metric, through a basic deformation, so that the projection of each $g$-minimal horizontal section becomes a harmonic immersion to the base?  

Express any solvability or obstruction condition in terms of known invariants: the O’Neill tensors $(A,T)$, the horizontal mean curvature, and the curvature of $(N,h)$. Is the only possibility the trivial case where $\phi$ has integrable horizontal distribution (so $(M,g)$ locally splits as a Riemannian product), or do non‑product examples exist?  

Why is it a "good" problem:  
This version asks sharply whether a **universal conformal gauge** exists aligning minimal and harmonic structures across the geometry of a harmonic morphism—an idea not directly resolved by standard classification theorems.  

1. **Profound Insight:** The question isolates a natural but untested structural principle: can the analytic notion of harmonicity of projections be made geometrically universal through a conformal rescaling? It thus probes the hidden symmetry between elliptic analytic conditions (harmonicity) and geometric minimality.  

2. **Bridge Between Fields:** It intertwines conformal geometry, Riemannian submersion theory (O’Neill tensors), and the analytic theory of harmonic morphisms. A solution would synthesize techniques from geometric analysis, submersion geometry, and potentially 2‑dimensional complex function theory on the base.  

3. **Extreme Simplicity:** Conceptually, the problem is easy to state—“Can one conformally adjust the metric so that minimal horizontal surfaces project harmonically?”—yet testing it requires deep understanding of curvature decompositions and conformal invariance of harmonic morphisms.  

4. **Open and Non-trivial:** Although formulas for $\tau(\phi|_\Sigma)$ are known, the simultaneous vanishing for *all* minimal sections under a single basic conformal rescaling has not been characterized or ruled out in general. Finding whether non‑product, nontrivial examples exist would either reveal a new rigidity phenomenon or a new class of conformally adapted harmonic morphisms—both significant and genuinely research‑level outcomes.